\documentclass{article}

\usepackage[preprint]{neurips_2026}

\usepackage[utf8]{inputenc}
\usepackage[T1]{fontenc}
\usepackage{hyperref}
\usepackage{url}
\usepackage{booktabs}
\usepackage{amsfonts}
\usepackage{amsmath}
\usepackage{amssymb}
\usepackage{nicefrac}
\usepackage{microtype}
\usepackage{xcolor}
\usepackage{graphicx}
\usepackage{enumitem}

\title{Physics-Informed Loss Formulations for Unrolled Power Flow Solvers}

\author{%
  Changhun Kim
}

\begin{document}

\maketitle

\section{Physics Loss Formulations}

\subsection{Default Physics Loss (MSE, No Normalization)}

The default physics loss is selected with:

\begin{verbatim}
--physics_loss_form mse
--physics_residual_norm none
--physics_final_weight 0.0
\end{verbatim}

All quantities are in \textbf{per-unit} (the dataset always normalizes by
$S_{\mathrm{base}}$, $U_{\mathrm{base}}$).

This combination activates a fast path in the implementation that directly
returns:

\begin{equation}
  \mathcal{L}_{\mathrm{phys}}
  =
  \sum_{k=0}^{K-1}
  \gamma^{K-1-k}
  \frac{1}{N}
  \sum_i
  \left[
    (\Delta P_i^{(k)})^2
    +
    (\Delta Q_i^{(k)})^2
  \right]
\end{equation}

where the slack/PV masking is applied \emph{before} the sum:
\begin{itemize}[noitemsep]
  \item $\Delta P_i$ is zeroed for the slack bus.
  \item $\Delta Q_i$ is zeroed for the slack bus and all PV buses.
\end{itemize}

Note: in this fast path the masked residuals are not filtered out; they
are set to zero and included in the mean, so $N$ is the total bus count.

%------------------------------------------------------------------------------
\section{Residual Normalization}

When \texttt{--physics\_residual\_norm} is not \texttt{none}, the code
first computes the per-bus apparent power magnitude from the scheduled
injections:

\begin{equation}
  |S_i^{\mathrm{set}}| = \sqrt{(P_i^{\mathrm{set}})^2 + (Q_i^{\mathrm{set}})^2}
\end{equation}

The masked residuals are then normalized, concatenated into a flat vector
$r^{(k)}$, and passed to the chosen loss function.  Only the \emph{active}
entries enter the mean: slack-bus $\Delta P$ and slack/PV-bus $\Delta Q$
entries are dropped, not zeroed.

\subsection{Setpoint Normalization (\texttt{--physics\_residual\_norm setpoint})}

Each bus is scaled by its own apparent-power setpoint, with a floor of
$\epsilon = 10^{-6}$ (p.u.) to guard against zero-injection buses:

\begin{equation}
  s_i = \max\!\bigl(|S_i^{\mathrm{set}}|,\;\epsilon\bigr)
\end{equation}

\begin{equation}
  \widetilde{\Delta P}_i^{(k)}
  =
  \frac{\Delta P_i^{(k)}}{s_i},
  \qquad
  \widetilde{\Delta Q}_i^{(k)}
  =
  \frac{\Delta Q_i^{(k)}}{s_i}
\end{equation}

This measures \emph{relative} mismatch: a 0.5 p.u.\ error on a 100 p.u.\
generator and on a 1 p.u.\ load contribute equally.

\subsection{Graph Normalization (\texttt{--physics\_residual\_norm graph})}

All buses in the same graph share a single scale equal to the maximum
apparent power in that graph (also floored at $\epsilon$):

\begin{equation}
  s_g = \max_{i \in g}\,|S_i^{\mathrm{set}}|
  \quad
  \text{(clipped to } \geq \epsilon\text{)}
\end{equation}

\begin{equation}
  \widetilde{\Delta P}_i^{(k)}
  =
  \frac{\Delta P_i^{(k)}}{s_g},
  \qquad
  \widetilde{\Delta Q}_i^{(k)}
  =
  \frac{\Delta Q_i^{(k)}}{s_g}
\end{equation}

The purpose is weaker than setpoint normalization: it puts all buses on the
same order of magnitude without trying to equalize small and large buses.

%------------------------------------------------------------------------------
\section{Selectable Loss Form}

After residual normalization (or in the non-\texttt{none} path), the
active residuals are collected into a flat vector:

\begin{equation}
  r^{(k)}
  =
  \bigl[
    \widetilde{\Delta P}^{(k)}_{\text{active}},\;
    \widetilde{\Delta Q}^{(k)}_{\text{active}}
  \bigr]
\end{equation}

and a penalty $\rho$ is applied element-wise.

\subsection{MSE (\texttt{--physics\_loss\_form mse})}

\begin{equation}
  \mathcal{L}_k = \frac{1}{M}\sum_j \bigl(r_j^{(k)}\bigr)^2
\end{equation}

\subsection{Huber (\texttt{--physics\_loss\_form huber})}

\begin{equation}
  \rho_\delta(r)
  =
  \begin{cases}
    \dfrac{1}{2}r^2, & |r| \le \delta \\[6pt]
    \delta\!\left(|r|-\dfrac{1}{2}\delta\right), & |r| > \delta
  \end{cases}
\end{equation}

The threshold $\delta$ is set with \texttt{--physics\_huber\_delta}
(default $1.0$).  Behaves like MSE near zero and like $\ell_1$ for large
residuals, so extreme buses dominate the gradient less.

\subsection{Log-Cosh (\texttt{--physics\_loss\_form logcosh})}

Implemented in a numerically stable form:

\begin{equation}
  \rho(r)
  = r + \log\!\bigl(1 + e^{-2r}\bigr) - \log 2
  = \log(\cosh(r))
\end{equation}

which approximates:

\begin{equation}
  \log(\cosh(r))
  \approx
  \begin{cases}
    \dfrac{1}{2}r^2 & \text{for small } |r| \\[4pt]
    |r| - \log 2   & \text{for large } |r|
  \end{cases}
\end{equation}

%------------------------------------------------------------------------------
\section{Final Residual Penalty}

The base loss sums discounted per-step terms:

\begin{equation}
  \mathcal{L}_{\mathrm{base}}
  =
  \sum_{k=0}^{K-1}
  \gamma^{K-1-k}
  \mathcal{L}_k
\end{equation}

When \texttt{--physics\_final\_weight}~$\lambda \neq 0$, the mismatch of
the \emph{final} output voltage is recomputed and appended:

\begin{equation}
  \mathcal{L}_{\mathrm{total}}
  =
  \mathcal{L}_{\mathrm{base}}
  \;+\;
  \lambda\,\mathcal{L}_{\mathrm{final}}
\end{equation}

where $\mathcal{L}_{\mathrm{final}}$ is computed with the same
normalization and loss form as the unrolled terms.  Importantly, it is
\emph{not} weighted by $\gamma$; it enters the sum with coefficient
$\lambda$ directly.

\paragraph{Reporting.}
Training metrics report raw (un-normalized) $\Delta P,\,\Delta Q$ in
per-unit.  Normalization affects only the gradient signal, not evaluation.

\end{document}
